% !PRINTED 496
\weberpage{496}

%<<CONT
the case in which the coefficients belong to a definite field
$\Omega$. We now make the assumption that this field $\Omega$ is an
\gesperrt{algebraic number field}, and accordingly consider expressions
$\varphi(x, y, z, \ldots)$ which are represented as a sum of terms of the
form
\[
\alpha x^r y^s z^t \ldots
\]
in which the exponents $r, s, t, \ldots$ are positive, or at least
non-negative, integers, while the coefficients $\alpha$ denote numbers in
$\Omega$. Such an expression
\weq{1}{\varphi(x, y, z, \ldots) = \sum \alpha x^r y^s z^t \ldots}
we call an \gesperrt{integral function in $\Omega$}. We always take the
expression to be so ordered and collected that the same combination of the
exponents $r, s, t, \ldots$ does not occur twice in it, and we call two such
expressions equal to one another only when they contain the same products
$x^r y^s z^t \ldots$, affected with the same coefficients. \gesperrt{An integral
function is set equal to zero when and only when all its coefficients are
zero.}

Several integral functions always yield integral functions again by
addition, subtraction and multiplication. By Bd.\,I, \S.\,143, I. one can
substitute for the variables $x, y, z, \ldots$ such rational numerical
values that one, or any number, of given integral functions in $\Omega$
different from zero receive non-vanishing numerical values (in $\Omega$).
From this it follows \gesperrt{that a product of several integral functions vanishes
only when one of its factors vanishes.}

The sum $r + s + t + \cdots$ of the exponents in a term of the expression
(1) is called the \gesperrt{degree of this term}, and the greatest value
which the degree of a term with non-vanishing coefficient in $\varphi$
assumes is called the \gesperrt{degree of the function $\varphi$}.

The degree of a product of two or more integral functions is equal to the
sum of the degrees of the individual factors.

For if in each of the factors one collects the sum of the terms of highest
degree into a homogeneous function, one obtains the terms of highest degree
of the product by multiplying all these homogeneous functions together.
